\section{Introduction}\label{sec:Intro}
For more than a century, magnetic core loss has been characterized by a single compact expression: the Steinmetz Equation (SE) and its variants \cite{1457110, 5955126}. The persistence of this formulation reflects its analytical value. The scaling relation $P_v=k f^\alpha B^\beta$ under sinusoidal drive conveys directly how loss depends on frequency and flux density, providing the quantitative insight on magnetic design. However, its limitation is that a three-parameter power law cannot precisely represent the nonlinear behavior of magnetic materials in a wide operating range, leading to structural errors and poor accuracy \cite{7370781}. 

To improve the accuracy, recent approaches have largely relinquished the explicit equation rather than refining it. Random forest \cite{10509368} and feedforward neural networks \cite{10502151} reproduce measured core loss to within a few percent errors by absorbing the loss nonlinearities into a learned mapping. However, they yield a numerical predictor rather than an analytical expression, and improved accuracy does not restore the scaling information that the SE conveys explicitly, which is insufficient for practical magnetic design. 

To bridge this trade-off, this paper proposes a Sparse Identification of Nonlinear Functions (SINF) framework that formulates equation discovery as a sparse selection problem over an expanded symbolic library. The dominant functions as well as their optimized coefficients are obtained by data-driven identification, forming the ultimate core loss model with simple expressions. In experiments, $7$ terms are retained out of $103$ candidates and the resulting explicit equation attains an $R^2$ of $0.9997$ and a MAPE of $1.56\%$ on Fair-Rite 95 data, achieving accuracy comparable to black-box models while preserving analytical interpretability.